% !PRINTED 22
\weberpage{22}
Mit Hilfe der Kovarianten haben wir die Form $f(x)$ auf die
\gesperrt{kanonische Form}
\[
\varphi(y) = y_1^3+y_2^3+y_3^3+6my_1y_2y_3
\]
reduziert.

Wir haben, wenn $r$ die Substitutionsdeterminante bedeutet:
\begin{align*}
P &= \frac{(2+m^3)m^3f^2+(2-5m^3)mr^2f\Delta+3m^2r^4\Delta^4-r^6J}{(1+8m^3)^2}, % [sic] printed \Delta^4; the later specialization uses \Delta^2
\\
Q &= \frac{(1+2m^3)f-6mr^2\Delta}{1+8m^3}, \\
R &= \frac{m^2f+r^2\Delta}{1+8m^3}
\end{align*}
gesetzt und erhielten $y_1^3$, $y_2^3$, $y_3^3$ als Wurzeln der kubischen
Gleichung
\[
u^3-Qu^2+Pu-R^3=0,
\]
und die Diskriminante dieser kubischen Gleichung ist
\[
D_1 = \frac{r^{18}K^2}{(1+8m^3)^6}.
\]

Drückt man diese Diskriminante durch $P$, $Q$, $R$ aus, so erhält
man $K^2$ rational durch $f$, $\Delta$, $J$ dargestellt. Wir haben an der
erwähnten Stelle der Algebra diesen Ausdruck nicht explizite angegeben. Jetzt müssen wir ihn aber bilden, wenn auch nur unter
der Voraussetzung $f=0$, d.\,h. nur für die Punkte der Kurve.
Es ist aber [nach Bd.\,I, § 50, (10)]
\[
D_1 = P^2Q^2+18PQR^3-4P^3-4Q^3R^3-27R^6.
\]
Darin hat man zu setzen:
\[
P = \frac{3m^2r^4\Delta^2-r^6J}{(1+8m^3)^2}, \quad Q = \frac{-6mr^2\Delta}{1+8m^3}, \quad R = \frac{r^2\Delta}{1+8m^3},
\]
und man erhält durch eine nicht schwierige Rechnung, wenn man
mit $S$ und $T$ die beiden \gesperrt{Invarianten} der Kurve dritter Ordnung
bezeichnet (Bd.\,II, § 108),
\[
(1+8m^3)^6D_1 = r^{18}(4J^3+108SJ\Delta^4-27T\Delta^6),
\]
und folglich
\weq{8}{K^2 = 4J^3+108SJ\Delta^4-27T\Delta^6.}
Diese Gleichung ist aber nicht identisch, sondern nur unter
der Voraussetzung $f=0$ befriedigt. Der vollständige Ausdruck
von $K^2$ durch $J$, $\Delta$, $f$ wird sehr viel komplizierter.
